\subsection{Product Perspective}

Dikaios Amanah is a web-based software platform designed to facilitate
structured interaction between verified Small and Medium Enterprises (SMEs),
retail investors, and system administrators.

The system follows a layered three-tier architecture consisting of a
presentation layer, an application layer, and a data layer. The separation
of these layers allows user interaction, application logic, security
controls, and persistent data management to be handled independently.

The presentation layer provides the web interface for Business Owners,
Investors, and Administrators. The application layer processes requests,
enforces business rules, manages authentication and authorization, and
coordinates the major platform services. The data layer manages operational
records, business information, investment information, and supporting
documents.

The proposed system is intended to operate as a technology-enabled
intermediary rather than as a financial institution. Investment transactions
within the academic MVP will be simulated for demonstration and testing
purposes.


\subsection{Product Functions}

The major functions of Dikaios Amanah are organized according to the
responsibilities of the system's primary users.

\subsubsection{User Registration and Authentication}

The system will provide account registration and authentication mechanisms
for users. Authentication will use secure password handling and token-based
session management. Access to system resources will be controlled according
to the user's assigned role.

\subsubsection{Business Owner Functions}

Business Owners will be able to:

\begin{itemize}
    \item Register an account.
    \item Complete identity and business verification.
    \item Manage business profile information.
    \item Submit supporting documents.
    \item Create investment campaigns after approval.
    \item Monitor campaign progress.
    \item Provide business progress information.
    \item Declare business profits.
    \item Communicate relevant information to investors.
\end{itemize}

A business must complete the required verification and administrative review
before its investment campaign becomes available to investors.


\subsubsection{Investor Functions}

Investors will be able to:

\begin{itemize}
    \item Register an account.
    \item Complete the required KYC verification process.
    \item Discover verified businesses and approved campaigns.
    \item Review business and campaign information.
    \item Review available risk information.
    \item Participate in simulated investment transactions.
    \item Monitor investments.
    \item View portfolio information.
    \item Track investment history.
    \item Monitor declared profits.
\end{itemize}


\subsubsection{Administrator Functions}

Administrators will be responsible for maintaining platform integrity and
managing operational activities. Administrative functions will include:

\begin{itemize}
    \item User verification.
    \item Business approval and rejection.
    \item Document validation.
    \item Risk assessment management.
    \item Campaign moderation.
    \item User management.
    \item Audit monitoring.
    \item Platform activity monitoring.
    \item Platform reporting.
\end{itemize}

Administrative actions will be recorded through the system's activity and
audit mechanisms.


\subsubsection{Investment Campaign Management}

The system will allow approved Business Owners to create and manage
investment campaigns. Campaign information will be presented to Investors
after the required administrative approval.

Investors will be able to review approved campaigns before participating.
The campaign lifecycle will support the proposed workflow from campaign
publication through investment participation and eventual completion.


\subsubsection{Document Management}

The system will allow Business Owners to submit supporting documents during
the verification process. Documents will be stored using the planned secure
cloud storage mechanism.

Cryptographic hashing will be used to support document integrity
verification and help identify unauthorized modification of stored
document
content.


\subsubsection{Trust and Audit Functions}

The system will maintain records of important platform events through a
Permissioned Trust and Audit Layer.

Critical events identified by the project proposal include:

\begin{itemize}
    \item Business verification.
    \item Contract approval.
    \item Investment confirmation.
    \item Profit declaration.
    \item Audit records.
    \item Document integrity verification.
\end{itemize}

The trust and audit mechanism will complement the primary relational
database rather than replace it.


\subsubsection{Notification and Activity Functions}

The system will provide notifications and activity information to help
users remain aware of relevant actions and changes within the platform.

Examples include verification status changes, campaign-related updates,
investment-related activities, and administrative actions.


\subsection{User Classes and Characteristics}

Dikaios Amanah contains three primary user classes.

\subsubsection{Business Owner}

A Business Owner represents an SME seeking investment through the platform.
Business Owners are expected to provide accurate business information and
supporting documentation. They will primarily interact with the platform
through the Business Owner Dashboard.

\subsubsection{Investor}

An Investor represents an individual seeking to participate in verified
Musharakah-based investment opportunities. Investors will use the platform
to review business information, evaluate available opportunities, manage
simulated investments, and monitor their portfolios.

\subsubsection{Administrator}

An Administrator is an authorized platform operator responsible for
verification, approval, monitoring, moderation, and other administrative
activities. Administrators have access to functions that are restricted
from Business Owners and Investors.


\subsection{Operating Environment}

The proposed system will operate as a web-based application.

The planned software environment consists of:

\begin{itemize}
    \item React.js for the frontend application.
    \item Node.js as the backend runtime environment.
    \item Express.js for backend and RESTful API services.
    \item PostgreSQL for relational data management.
    \item Firebase Storage for secure document storage.
    \item JWT for authentication.
    \item bcrypt for password hashing.
    \item SHA-256 for document integrity verification.
    \item Git and GitHub for version control.
\end{itemize}

Users will access the system through modern web browsers using an internet
connection. Communication between the frontend and backend services will
use HTTPS.


\subsection{Design Constraints}

The following constraints are derived from the defined project scope and
academic MVP objectives:

\begin{itemize}
    \item The project must remain within the defined academic MVP scope.
    \item Investment transactions will be simulated rather than connected to
    real financial payment systems.
    \item Live banking integration is excluded from the MVP.
    \item Real payment gateway integration is excluded from the MVP.
    \item Production-level BSEC regulatory compliance implementation is
    excluded from the MVP.
    \item AI-based fraud detection is excluded from the MVP.
    \item AI-powered risk prediction is excluded from the MVP.
    \item Production Hyperledger Fabric deployment is excluded from the MVP.
    \item Automated smart contract execution is excluded from the MVP.
    \item Native Android and iOS applications are excluded from the MVP.
\end{itemize}

These limitations are intended to keep the system technically achievable
within the project's academic development scope.


\subsection{Assumptions and Dependencies}

The system development is based on several assumptions established in the
project proposal:

\begin{itemize}
    \item Users have basic digital literacy.
    \item Users have access to an internet connection.
    \item Businesses provide authentic documentation.
    \item Administrators perform verification activities responsibly.
    \item The system is initially developed for academic demonstration
    purposes.
    \item Investment transactions used during testing are simulated.
\end{itemize}

The system also depends on the availability and proper operation of its
planned software and infrastructure components, including the relational
database, document storage service, authentication mechanisms, and
deployment environment.